\chapter{Vector Bundles}\label{vector-bundles}
\chapterstatus{complete}

\section{Introduction}\label{vector-bundles:section-introduction}

A vector bundle is a family of vector spaces $E_p$, one for each point $p$ of a manifold $M$, glued together smoothly. The model is the tangent bundle: at each point of a
surface in $\RR^3$ there is a tangent plane, and as the point moves the plane moves with it. Locally such a family looks like a product $U \times \mathbb{K}^r$, a
\emph{local trivialisation}; globally it may be twisted, as the M\"obius band is a twisted version of the cylinder. The twisting is recorded by \emph{transition functions}:
where two trivialisations overlap, they differ by a smoothly varying invertible matrix.

Vector bundles appear everywhere in this book. Tangent and cotangent bundles and the bundles of differential forms are the setting of analysis on manifolds; sections of
bundles are ``twisted functions'', which is how the line bundles of complex and algebraic geometry encode divisors and linear systems; and the amount of twisting is measured
by characteristic classes, the first of which, the first Chern class of a complex line bundle, is constructed at the end of this chapter.

The chapter goes as follows.
\begin{enumerate}
\item Definitions, examples, and the description of bundles by cocycles of transition functions (Section~\ref{vector-bundles:section-definition}), with the M\"obius bundle,
the tangent bundle and the tautological bundles worked out.
\item Sections and frames, and the dictionary between vector bundles and locally free sheaves.
\item The linear algebra of bundles: direct sums, tensor products, duals, $\Hom$ bundles, exterior powers, pullbacks; metrics, subbundles, kernels and normal bundles
(Section~\ref{vector-bundles:section-operations}).
\item Complex line bundles: their classification by $H^2(M; \ZZ)$, and an elementary classification on the sphere by the winding number of a transition function
(Section~\ref{vector-bundles:section-line-bundles}).
\end{enumerate}
References: \cite{lee-smooth-manifolds}, \cite{milnor-stasheff}, \cite[Chapter 0]{griffiths-harris}.

\noindent\textbf{Prerequisites.} Chapter~\ref{smooth-manifolds} (Smooth Manifolds), Chapter~\ref{multilinear-algebra} (Multilinear Algebra) and Chapter~\ref{sheaves} (Sheaves on
Topological Spaces).

\noindent\textbf{Conventions.} $\mathbb{K} = \RR$ or $\CC$, and $M$ is a smooth manifold. Vectors of $\mathbb{K}^r$ are columns, and $e_1, \ldots, e_r$ is the standard basis.
$\GL_r(\mathbb{K})$ is the group of invertible $r \times r$ matrices, an open subset of $M_r(\mathbb{K})$.

\section{Vector bundles and cocycles}\label{vector-bundles:section-definition}

\begin{definition}\label{vector-bundles:definition-vector-bundle}
A \emph{vector bundle} of rank $r$ over $M$ is a smooth map $\pi \colon E \to M$ together with $\mathbb{K}$-vector space structures on the fibres $E_p = \pi^{-1}(p)$, such that every point of $M$ has an open
neighbourhood $U$ and a diffeomorphism $\Phi \colon \pi^{-1}(U) \to U \times \mathbb{K}^r$ (a \emph{local trivialisation}) with $\mathrm{pr}_1 \circ \Phi = \pi$, linear on each fibre. A \emph{morphism} of vector
bundles over $M$ is a smooth $f \colon E \to F$ with $\pi_F \circ f = \pi_E$, linear on fibres; \emph{line bundle} means rank $1$.
\end{definition}

Unwinding the definition: $E$ is a manifold, $\pi$ maps it onto $M$, each fibre $E_p$ is a vector space of dimension $r$, and over a neighbourhood of each point the manifold $E$
is identified with $U \times \mathbb{K}^r$ in such a way that the fibre $E_p$ goes to $\{p\} \times \mathbb{K}^r$ by a linear isomorphism $\Phi_p \colon E_p \to \mathbb{K}^r$.
The rank is part of the data; one could allow it to vary, but it is locally constant (because of the local trivialisations), so on a connected manifold it is constant. An
\emph{isomorphism} is a morphism with an inverse morphism, and $E$ is \emph{trivial} if it is isomorphic to the \emph{trivial bundle} $M \times \mathbb{K}^r$ with the projection
to $M$. A complex vector bundle of rank $r$ is also a real one of rank $2r$, by forgetting the complex structure.

\begin{lemma}\label{vector-bundles:lemma-bijective-morphism}
A morphism of vector bundles over $M$ that is bijective on every fibre is an isomorphism.
\end{lemma}

\begin{proof}
The inverse map is linear on fibres; we must show it is smooth. This is local, so we may assume both bundles are trivial, $U \times \mathbb{K}^r$, and the morphism is
$(p, v) \mapsto (p, A(p)v)$ for a smooth map $A \colon U \to M_r(\mathbb{K})$ (its entries are the components of the images of the constant sections $e_j$). By hypothesis $A(p)$ is
invertible, and $A(p)^{-1}$ is smooth in $p$ by Cramer's rule, since the entries of the inverse are polynomials in the entries of $A$ divided by $\det A \neq 0$. So the inverse
$(p, w) \mapsto (p, A(p)^{-1}w)$ is smooth.
\end{proof}

\begin{example}\label{vector-bundles:example-bundles}
\begin{enumerate}
\item \emph{The trivial bundle} $M \times \mathbb{K}^r$.
\item \emph{The tangent bundle} $TM$ (Definition~\ref{smooth-manifolds:definition-tangent-bundle}). A chart $\varphi = (x^1, \ldots, x^n)$ on $U$ gives the local trivialisation
$T_pM \to \RR^n$, $v \mapsto (dx^1(v), \ldots, dx^n(v))$, sending $\sum_iv^i\partial/\partial x^i$ to $(v^1, \ldots, v^n)$. Its dual, the \emph{cotangent bundle} $T^*M$, has fibres
$T_p^*M = (T_pM)^*$.
\item \emph{The M\"obius bundle} over $S^1 = \RR/\ZZ$ is the quotient of $\RR \times \RR$ by $(x, t) \sim (x + 1, -t)$, with the projection $[x, t] \mapsto [x]$. Equivalently,
it is obtained from $[0, 1] \times \RR$ by gluing $(0, t)$ to $(1, -t)$. Over each interval $\{[x] : a < x < a + 1\}$ it is trivialised by $[x, t] \mapsto ([x], t)$ with
$a < x < a + 1$.
\item \emph{The tautological line bundles.} On $\RR\PP^n$ and on $\CC\PP^n$, $\mathcal{O}(-1) = \{(\ell, v) : v \in \ell\} \subseteq \mathbb{K}\PP^n \times \mathbb{K}^{n+1}$, whose fibre over a
line $\ell$ is the line itself (Exercise~\ref{vector-bundles:exercise-tautological}).
\item \emph{The normal bundle} $N_{S/M} = (TM|_S)/TS$ of a submanifold $S \subseteq M$ (Proposition~\ref{vector-bundles:proposition-constant-rank}); for $S \subseteq \RR^N$ it can be
realised as $\{(p, v) \in S \times \RR^N : v \perp T_pS\}$.
\end{enumerate}
\end{example}

\begin{proposition}\label{vector-bundles:proposition-lie-group-parallelisable}
The tangent bundle of a Lie group $G$ is trivial: $G \times \mathfrak{g} \to TG$, $(g, v) \mapsto d(L_g)_e(v)$, is an isomorphism of vector bundles. In particular $S^1$, $S^3 \cong \mathrm{SU}(2)$ and the tori $T^n$ have trivial tangent
bundles.
\end{proposition}

\begin{proof}
The map is smooth, as the differential of multiplication restricted to $G \times T_eG$, it covers the identity of $G$, and it is a linear isomorphism on each fibre since $L_g$ is a diffeomorphism. By
Lemma~\ref{vector-bundles:lemma-bijective-morphism} it is an isomorphism. For $S^3$ note that
$\mathrm{SU}(2) = \{\begin{pmatrix}a & -\bar b\\ b & \bar a\end{pmatrix} : |a|^2 + |b|^2 = 1\}$ is diffeomorphic to $S^3 \subseteq \CC^2$.
\end{proof}

\begin{counterexample}\label{vector-bundles:counterexample-mobius}
The M\"obius bundle $E$ over $S^1$ (Example~\ref{vector-bundles:example-bundles}) is not trivial. A nowhere vanishing section would correspond to a continuous $s \colon [0, 1] \to \RR \setminus \{0\}$ with $s(1) = -s(0)$, and by the
intermediate value theorem $s$ would vanish somewhere; and a real line bundle with a nowhere vanishing section $s$ is trivial, by $(p, t) \mapsto ts(p)$. The tautological line bundle of $\RR\PP^1$ is isomorphic to $E$: parametrising
$\RR\PP^1$ by $\theta \in [0, \pi]$, with $\theta$ corresponding to the line through $u(\theta) = (\cos\theta, \sin\theta)$, the section $u$ trivialises the bundle over $[0, \pi]$ and $u(\pi) = -u(0)$, which is the gluing of $E$. Hence the
tautological line bundle of $\RR\PP^n$, which restricts to it on $\RR\PP^1 \subseteq \RR\PP^n$, is nontrivial as well.
\end{counterexample}

Two local trivialisations $\Phi_i$ and $\Phi_j$ of the same bundle can be compared on the overlap of their domains: $\Phi_i \circ \Phi_j^{-1}$ maps $(p, v)$ to
$(p, g_{ij}(p)v)$, where $g_{ij}(p) = \Phi_{i,p} \circ \Phi_{j,p}^{-1}$ is the invertible matrix that converts the coordinates of a vector of $E_p$ in the trivialisation $j$ into
its coordinates in the trivialisation $i$. These matrices are the transition functions, and they determine the bundle.

\begin{proposition}\label{vector-bundles:proposition-cocycles}
Let $(U_i)$ be an open cover of $M$.
\begin{enumerate}
\item A vector bundle with trivialisations $\Phi_i$ over $U_i$ determines \emph{transition functions} $g_{ij} \colon U_i \cap U_j \to \GL_r(\mathbb{K})$, smooth, with
$\Phi_i\Phi_j^{-1}(p, v) = (p, g_{ij}(p)v)$; they satisfy the \emph{cocycle condition} $g_{ij}g_{jk} = g_{ik}$ on triple intersections, and $g_{ii} = 1$.
\item Conversely, every such family $(g_{ij})$ arises from a vector bundle, unique up to isomorphism, namely $E = (\bigsqcup_i U_i \times \mathbb{K}^r)/\sim$ with $(p, v)_j \sim (p, g_{ij}(p)v)_i$.
\item Two cocycles $(g_{ij})$, $(g'_{ij})$ define isomorphic bundles if and only if there are smooth $h_i \colon U_i \to \GL_r(\mathbb{K})$ with $g'_{ij} = h_ig_{ij}h_j^{-1}$. Hence isomorphism classes of rank
$r$ bundles correspond to $\check{H}^1(M, \underline{\GL_r})$, \v{C}ech cohomology with values in the sheaf of smooth $\GL_r(\mathbb{K})$-valued functions
(Remark~\ref{sheaf-cohomology:remark-cech-limit}).
\end{enumerate}
\end{proposition}

\begin{proof}
(1) The map $\Phi_i\Phi_j^{-1}$ is a diffeomorphism of $(U_i \cap U_j) \times \mathbb{K}^r$ that commutes with the projection to $U_i \cap U_j$ and is linear on fibres, so it has
the form $(p, v) \mapsto (p, g_{ij}(p)v)$ with $g_{ij}(p)$ invertible. The entries of $g_{ij}$ are smooth: the $k$-th column of $g_{ij}(p)$ is the second component of
$\Phi_i\Phi_j^{-1}(p, e_k)$, a smooth function of $p$. On $U_i \cap U_j \cap U_k$, $\Phi_i\Phi_k^{-1} = (\Phi_i\Phi_j^{-1})(\Phi_j\Phi_k^{-1})$, which is $g_{ik} = g_{ij}g_{jk}$, and $g_{ii} = 1$.
In particular $g_{ji} = g_{ij}^{-1}$.

(2) The relation $\sim$ is reflexive because $g_{ii} = 1$, symmetric because $g_{ji} = g_{ij}^{-1}$, and transitive because of the cocycle condition. Let $E$ be the quotient,
$\pi[p, v]_i = p$, and give $E_p$ the vector space structure transported from $\mathbb{K}^r$ by $v \mapsto [p, v]_i$ for any $i$ with $p \in U_i$; this does not depend on $i$, as the
identifications differ by the linear maps $g_{ij}(p)$. The maps $\Phi_i \colon \pi^{-1}(U_i) \to U_i \times \mathbb{K}^r$, $[p, v]_i \mapsto (p, v)$, are bijective, and their transition maps
$(p, v) \mapsto (p, g_{ij}(p)v)$ are diffeomorphisms of the open sets $(U_i \cap U_j) \times \mathbb{K}^r$. Composing with charts of $M$, the $\Phi_i$ form an atlas of $E$ (Hausdorff and second
countable because $M$ is and the cover can be taken countable), for which $\pi$ is smooth and the $\Phi_i$ are local trivialisations with transition functions $g_{ij}$. If a bundle
$E^{\prime}$ has trivialisations $\Phi^{\prime}_i$ with the same transition functions, then $\Phi_i^{\prime-1}\Phi_i$ on $\pi^{-1}(U_i)$ agree on overlaps and glue to an isomorphism $E \to E^{\prime}$.

(3) An isomorphism $f \colon E \to E^{\prime}$ is given in the trivialisations by $\Phi^{\prime}_if\Phi_i^{-1}(p, v) = (p, h_i(p)v)$ with $h_i$ smooth and invertible (as in (1)). On
$U_i \cap U_j$, computing $\Phi^{\prime}_if\Phi_j^{-1}$ in two ways gives $h_ig_{ij} = g^{\prime}_{ij}h_j$. Conversely, given such $h_i$, the maps $\Phi_i^{\prime-1}h_i\Phi_i$ agree on overlaps and define
an isomorphism, which is smooth with smooth inverse. The relation $g^{\prime}_{ij} = h_ig_{ij}h_j^{-1}$ is exactly the statement that the two \v{C}ech cocycles are cohomologous.
\end{proof}

\begin{example}\label{vector-bundles:example-cocycles}
\begin{enumerate}
\item \emph{The tangent bundle.} For the trivialisations of $TM$ given by charts $\varphi$ and $\psi$ (Example~\ref{vector-bundles:example-bundles}(2)), a vector with components
$\xi$ in the chart $\varphi$ has components $D\tau(\varphi(p))\xi$ in the chart $\psi$, where $\tau = \psi \circ \varphi^{-1}$ (Proposition~\ref{smooth-manifolds:proposition-tangent-bundle}).
So the transition functions of $TM$ are the Jacobian matrices of the transition maps of the atlas, and the cocycle condition is the chain rule.
\item \emph{The sphere.} For the two stereographic charts of $S^n$ (Example~\ref{general-topology:example-stereographic}), the transition map is the inversion
$\tau(y) = y/|y|^2$ of $\RR^n \setminus \{0\}$, whose Jacobian matrix is
\[
D\tau(y) = \frac{1}{|y|^2}\Bigl(I - \frac{2yy^T}{|y|^2}\Bigr),
\]
the reflection in the hyperplane $y^\perp$ divided by $|y|^2$. It has determinant $-|y|^{-2n}$: the two charts induce opposite orientations. This single matrix-valued function
on $\RR^n \setminus \{0\}$ describes $TS^n$ completely.
\item \emph{The M\"obius bundle.} With the two trivialisations of Example~\ref{vector-bundles:example-bundles}(3) over the intervals $0 < x < 1$ and $\frac12 < x < \frac32$, the overlap
consists of two arcs; on one the transition function is $1$ and on the other it is $-1$. A cocycle with values in $\{\pm1\} \subseteq \GL_1(\RR)$ describes a bundle with locally constant
transition functions, a \emph{flat} bundle (Chapter~\ref{connections}).
\item \emph{The tautological line bundle of $\CC\PP^1$.} Over $U_0 = \{[1 : z]\}$ and $U_1 = \{[w : 1]\}$ the line $\ell$ has the basis vectors $s_0 = (1, z)$ and $s_1 = (w, 1)$, which
trivialise $\mathcal{O}(-1)$. On $U_0 \cap U_1$, where $w = 1/z$, $s_1 = z^{-1}s_0$. So a vector $v = a\,s_0 = b\,s_1$ has coordinates related by $a = z^{-1}b$: the transition function
is $g_{01} = z^{-1}$, a function $\CC^\times \to \GL_1(\CC) = \CC^\times$.
\end{enumerate}
\end{example}

\begin{definition}\label{vector-bundles:definition-sections}
A \emph{section} of $E$ over $U \subseteq M$ open is a smooth $s \colon U \to E$ with $\pi \circ s = \id$; they form a $C^\infty(U)$-module $\Gamma(U, E)$, and $U \mapsto \Gamma(U, E)$ is a sheaf
$\mathcal{E}$ of $C^\infty_M$-modules. A \emph{frame} over $U$ is a tuple of sections that is a basis of each fibre.
\end{definition}

A section picks a vector in each fibre, smoothly. In a trivialisation $\Phi$ over $U$ a section is the same as a smooth function $U \to \mathbb{K}^r$, namely $\Phi \circ s$; so
sections are ``functions twisted by the bundle''. The sections of the trivial line bundle are the functions, the sections of $TM$ are the vector fields, and those of $T^*M$ are the
$1$-forms. Every bundle has the zero section, and many others: if $v \in E_p$, a bump function times a local section through $v$ gives a global section with value $v$ at $p$.

\begin{example}\label{vector-bundles:example-sections}
\begin{enumerate}
\item \emph{Sections of the M\"obius bundle} are the smooth functions $f \colon \RR \to \RR$ with $f(x + 1) = -f(x)$, such as $\sin(\pi x)$, which vanishes exactly at the integers.
Every section vanishes somewhere (Counterexample~\ref{vector-bundles:counterexample-mobius}).
\item \emph{Sections of a bundle given by a cocycle} are families of functions $s_i \colon U_i \to \mathbb{K}^r$ with $s_i = g_{ij}s_j$ on $U_i \cap U_j$. For $\mathcal{O}(-1)$ on
$\CC\PP^1$, with Example~\ref{vector-bundles:example-cocycles}(4), a section is a pair of functions $a$ on $U_0$ and $b$ on $U_1$ with $a = z^{-1}b$; holomorphic sections would be
power series $a(z)$ and $b(w)$ with $a(z) = z^{-1}b(1/z)$, and comparing the powers of $z$ shows that only $a = b = 0$ is possible. So $\mathcal{O}(-1)$ has no nonzero holomorphic
section, although it has many smooth ones (Chapter~\ref{holomorphic-bundles}).
\end{enumerate}
\end{example}

\begin{proposition}\label{vector-bundles:proposition-sheaf-of-sections}
$\mathcal{E}$ is locally free of rank $r$ (Definition~\ref{sheaves:definition-module-operations}): a trivialisation over $U$ is the same as a frame over $U$, and gives $\mathcal{E}|_U \cong (C^\infty_U)^r$.
Conversely every locally free $C^\infty_M$-module of rank $r$ is the sheaf of sections of a vector bundle, unique up to isomorphism, and this correspondence is an equivalence between vector bundles and
locally free sheaves.
\end{proposition}

\begin{proof}
A trivialisation $\Phi$ gives the frame $s_i(p) = \Phi^{-1}(p, e_i)$, and conversely a frame defines $\Phi(v) = (p, (v^1, \ldots, v^r))$ where $v = \sum v^is_i(p)$; smoothness in both directions follows because
the coefficients are smooth functions (they are the components in a trivialisation). Given a locally free sheaf $\mathcal{F}$, choose isomorphisms $\mathcal{F}|_{U_i} \cong (C^\infty_{U_i})^r$; on overlaps
they differ by automorphisms of $(C^\infty)^r$, that is, by matrices $g_{ij}$ of smooth functions with smooth inverse, which form a cocycle. The bundle of
Proposition~\ref{vector-bundles:proposition-cocycles}(2) has $\mathcal{F}$ as its sheaf of sections. Morphisms correspond in the same way, since a morphism of bundles is determined by its action on sections.
\end{proof}

\section{Operations on bundles}\label{vector-bundles:section-operations}

Every construction of linear algebra that is applied to vector spaces can be applied fibre by fibre to vector bundles. The principle is always the same: apply the construction to
the transition functions, and use Proposition~\ref{vector-bundles:proposition-cocycles}(2).

\begin{proposition}\label{vector-bundles:proposition-operations}
Let $E$, $F$ be vector bundles over $M$ of ranks $r$, $s$, with cocycles $(g_{ij})$, $(h_{ij})$, and let $f \colon N \to M$ be smooth.
\begin{enumerate}
\item There are vector bundles $E \oplus F$, $E \otimes F$, $E^\vee$, $\Hom(E, F)$, $\Lambda^kE$ and $\mathrm{Sym}^kE$ whose fibres are the corresponding constructions on the fibres, with cocycles
$g_{ij} \oplus h_{ij}$, $g_{ij} \otimes h_{ij}$, $(g_{ij}^T)^{-1}$, and so on. Their sheaves of sections are the corresponding operations on locally free sheaves.
\item The \emph{pullback} $f^*E$, with fibre $(f^*E)_q = E_{f(q)}$ and cocycle $g_{ij} \circ f$ for the cover $(f^{-1}U_i)$, is a vector bundle on $N$; $f^*$ is functorial and
$(g \circ f)^* \cong f^*g^*$.
\item $\det E = \Lambda^rE$ is a line bundle with cocycle $\det g_{ij}$.
\end{enumerate}
\end{proposition}

Explicitly, the cocycles are: for $E \oplus F$ the block matrices $\bigl(\begin{smallmatrix} g_{ij} & 0\\ 0 & h_{ij}\end{smallmatrix}\bigr)$; for $E \otimes F$ the Kronecker products
$g_{ij} \otimes h_{ij}$, which act on $v \otimes w$ by $g_{ij}v \otimes h_{ij}w$; for the dual $E^\vee$, $(g_{ij}^T)^{-1}$, because a linear form with coordinates $\xi$ (a row vector
$\xi^T$) must satisfy $\xi_i^Tv_i = \xi_j^Tv_j$ when $v_i = g_{ij}v_j$; for $\Hom(E, F)$, the map $A \mapsto h_{ij}Ag_{ij}^{-1}$ on matrices; for $\Lambda^kE$ the matrices of $k \times k$
minors of $g_{ij}$; and for $\det E$ the determinants.

\begin{proof}
In each case the fibrewise construction is a functor on finite-dimensional vector spaces and isomorphisms, and the transition matrices of the new bundle are obtained from the
$g_{ij}$ and $h_{ij}$ by applying this functor. Their entries are polynomials in the entries of $g_{ij}$, $h_{ij}$, possibly divided by $\det g_{ij}$, so they are smooth; and functoriality
turns the cocycle conditions for $(g_{ij})$ and $(h_{ij})$ into the cocycle condition for the new family. By Proposition~\ref{vector-bundles:proposition-cocycles}(2) there is a bundle with
these transition functions, and its fibre over $p$ is canonically the construction applied to $E_p$ and $F_p$, since the identifications with $\mathbb{K}^r$ and $\mathbb{K}^s$ over different
$U_i$ differ exactly by the transition matrices. The sheaf of sections is computed locally, where everything is trivial. For the pullback, the functions $g_{ij} \circ f$ are smooth and
satisfy the cocycle condition on the cover $(f^{-1}U_i)$, and $(f^*E)_q = E_{f(q)}$ canonically; functoriality holds because the cocycle of $(g \circ f)^*E$ is that of $f^*(g^*E)$.
For the determinant, $\Lambda^r$ of an $r \times r$ matrix is multiplication by its determinant.
\end{proof}

\begin{example}\label{vector-bundles:example-operations}
\begin{enumerate}
\item The cotangent bundle $T^*M$ has transition functions $(D\tau^T)^{-1}$, the inverse transposes of the Jacobian matrices; the bundle $\Lambda^kT^*M$ of $k$-forms has the $k \times k$
minors of these, and the bundle $\Lambda^nT^*M$ of top degree forms has transition functions $\det(D\tau)^{-1}$. It is trivial when the atlas can be chosen with $\det D\tau > 0$,
which is the orientability of $M$ (Chapter~\ref{differential-forms}).
\item The line bundles $\mathcal{O}(k)$ on $\CC\PP^1$: $\mathcal{O}(1) = \mathcal{O}(-1)^\vee$ has transition function $z$, and $\mathcal{O}(k) = \mathcal{O}(1)^{\otimes k}$ for $k \geq 0$,
$\mathcal{O}(k) = \mathcal{O}(-1)^{\otimes(-k)}$ for $k < 0$, has transition function $z^k$ with respect to the covering $U_0$, $U_1$ (for line bundles, tensor products multiply the
transition functions, and duals invert them).
\item \emph{Restrictions and constant maps.} For the inclusion $\iota \colon S \to M$ of a submanifold, $\iota^*E = E|_S$ is the restriction. For a constant map $f \colon N \to M$ with value $p$,
$f^*E = N \times E_p$ is trivial.
\end{enumerate}
\end{example}

\begin{proposition}\label{vector-bundles:proposition-pullback-fibre-product}
Let $\pi \colon E \to M$ be a vector bundle and $f \colon N \to M$ smooth. The fibre product
\[
N \times_ME = \{(q, e) \in N \times E : f(q) = \pi(e)\}
\]
is a submanifold of $N \times E$ (Proposition~\ref{smooth-manifolds:proposition-fibre-product}), and the projection to $N$, with the vector space structures of the
fibres $E_{f(q)}$, is a vector bundle isomorphic to the pullback $f^*E$ of Proposition~\ref{vector-bundles:proposition-operations}(2). It has the universal
property of the fibre product: for a manifold $L$, the smooth maps $L \to f^*E$ correspond to pairs of smooth maps $a \colon L \to N$ and $b \colon L \to E$ with
$f \circ a = \pi \circ b$. In particular the sections of $f^*E$ over $N$ correspond to the smooth maps $s \colon N \to E$ with $\pi \circ s = f$, and a section $t$ of
$E$ pulls back to the section $f^*t = t \circ f$.
\end{proposition}

\begin{proof}
A vector bundle projection is a submersion, since in a local trivialisation it is the projection $U \times \mathbb{K}^r \to U$, so
Proposition~\ref{smooth-manifolds:proposition-fibre-product} applies and gives the submanifold and the universal property. If $\Phi \colon \pi^{-1}(U) \to U \times \mathbb{K}^r$
is a local trivialisation of $E$, then $(q, e) \mapsto (q, \mathrm{pr}_2\Phi(e))$ is a diffeomorphism from the part of $N \times_ME$ over $f^{-1}(U)$ onto
$f^{-1}(U) \times \mathbb{K}^r$, linear on fibres; so the projection to $N$ is a vector bundle, and its transition functions for these trivialisations are the
$g_{ij} \circ f$, which are those of $f^*E$. Two vector bundles with the same cocycle are isomorphic (Proposition~\ref{vector-bundles:proposition-cocycles}).
\end{proof}

So a section of $f^*E$ is ``a section of $E$ along $f$'': for example, for a curve $\gamma$ in $M$, a section of $\gamma^*TM$ is a vector field along the curve, such as its
velocity, and the differential $df$ of a smooth map is a section of $T^*N \otimes f^*TM$.

\begin{proposition}\label{vector-bundles:proposition-metrics}
Every vector bundle over a manifold admits a metric: a smooth family of inner products (Hermitian inner products in the complex case) on the fibres. Consequently:
\begin{enumerate}
\item every subbundle has a complement, and every short exact sequence $0 \to E' \to E \to E'' \to 0$ of vector bundles splits;
\item $E \cong E^\vee$ for real bundles, and $E^\vee \cong \overline{E}$ for complex bundles with Hermitian metric.
\end{enumerate}
\end{proposition}

Here a \emph{subbundle} of $E$ is a subset $E^{\prime} \subseteq E$ such that every point has a neighbourhood with a frame $s_1, \ldots, s_r$ of $E$ for which $s_1, \ldots, s_k$ is a frame of
$E^{\prime}$; then $E^{\prime}$ is a vector bundle of rank $k$, and the quotient $E/E^{\prime}$, with fibres $E_p/E^{\prime}_p$, is a vector bundle with the local frames given by the images of
$s_{k+1}, \ldots, s_r$. The \emph{conjugate} $\overline{E}$ of a complex bundle is the same real bundle with the conjugate complex structure.

\begin{proof}
Choose trivialisations over a cover $(U_i)$ and a subordinate partition of unity $(\psi_i)$ (Theorem~\ref{smooth-manifolds:theorem-partitions-of-unity}); the sum $\sum_i\psi_i\langle -, - \rangle_i$ of the
pulled-back standard inner products is a metric, since a convex combination of inner products is an inner product. Given a subbundle $E' \subseteq E$, the orthogonal complement $E'^\perp$ is a subbundle (in a
local frame adapted to $E'$, apply Gram--Schmidt, Theorem~\ref{linear-algebra:theorem-gram-schmidt}, which depends smoothly on the frame) and $E = E' \oplus E'^\perp$; a surjection $E \to E''$ then splits
by inverting $E'^\perp \to E''$. Statement (2) is the fibrewise isomorphism $v \mapsto \langle -, v \rangle$, which is smooth.
\end{proof}

With a metric one can always choose \emph{orthonormal} local frames, by Gram--Schmidt; then the transition functions are orthogonal (unitary) matrices. So every real bundle can
be described by a cocycle with values in $\mathrm{O}(r)$, and every complex one by a cocycle with values in $\mathrm{U}(r)$: the structure group can be \emph{reduced} to a compact group.

\begin{proposition}\label{vector-bundles:proposition-frames}
\begin{enumerate}
\item A vector bundle $E$ of rank $r$ is trivial if and only if it has $r$ sections $s_1, \ldots, s_r$ that are linearly independent at every point (a \emph{global frame}).
\item For a real line bundle $L$, $L \otimes L$ is trivial. For a complex line bundle $L$, $L \otimes \overline{L}$ is trivial, so $\overline{L} \cong L^\vee$.
\end{enumerate}
In particular the isomorphism classes of real line bundles form a group under $\otimes$ in which every element has order at most $2$.
\end{proposition}

\begin{proof}
(1) A frame gives the morphism $M \times \mathbb{K}^r \to E$, $(p, a) \mapsto \sum_ia_is_i(p)$, which is bijective on each fibre, hence an isomorphism
(Lemma~\ref{vector-bundles:lemma-bijective-morphism}); conversely the standard basis vectors are a frame of the trivial bundle. (2) Choose a metric (Proposition~\ref{vector-bundles:proposition-metrics}).
In the real case the map $L \otimes L \to M \times \RR$, $v \otimes w \mapsto \langle v, w\rangle$, is fibrewise a nonzero linear map between lines, hence an isomorphism; smoothness is checked in a local frame. In the
complex case use $v \otimes \bar w \mapsto \langle v, w\rangle$ on $L \otimes \overline{L}$. The group structure on line bundles is Proposition~\ref{vector-bundles:proposition-operations} and the uniqueness of
tensor products up to isomorphism, with the trivial bundle as unit and $L^\vee$ as inverse, since $L \otimes L^\vee \cong \End(L)$ has the nowhere vanishing section $\id$.
\end{proof}

\begin{example}\label{vector-bundles:example-mobius-square}
For the M\"obius bundle $E$ over $S^1$, $E \otimes E$ is trivial by Proposition~\ref{vector-bundles:proposition-frames}, although $E$ is not (Counterexample~\ref{vector-bundles:counterexample-mobius}). Directly: $E \otimes E$ is glued from
$[0, 1] \times \RR$ by $(0, t) \sim (1, (-1)^2t) = (1, t)$, which is the trivial bundle. Likewise $E \oplus E$ is trivial: it is glued by $-\id_{\RR^2}$, the rotation by $\pi$, and the rotations by $\pi s$, $0 \leq s \leq 1$,
connect it to the identity. Concretely, the sections $x \mapsto (\cos\pi x, \sin\pi x)$ and $x \mapsto (-\sin\pi x, \cos\pi x)$ of $E \oplus E$ (functions $\RR \to \RR^2$ with
$f(x + 1) = -f(x)$) form a global frame. The same argument shows that a real bundle glued from $[0, 1] \times \RR^r$ by a matrix $A$ is trivial when $A$ can be joined to the identity by a path in $\GL_r(\RR)$: if $A_t$ is such a path, constant near its ends, with $A_0 = \id$ and $A_1 = A$,
the sections $t \mapsto A_te_i$ satisfy the gluing condition and form a global frame.
\end{example}

\begin{proposition}\label{vector-bundles:proposition-constant-rank}
Let $f \colon E \to F$ be a morphism of vector bundles over $M$ whose rank on the fibres is constant. Then $\ker f \subseteq E$ and $\operatorname{im}f \subseteq F$ are subbundles, and $\operatorname{coker}f = F/\operatorname{im}f$ is a vector
bundle. For a submanifold $S \subseteq M$ the normal bundle $N_{S/M} = (TM|_S)/TS$ is a vector bundle, and if $S = \phi^{-1}(q)$ for a regular value $q$ of a smooth $\phi \colon M \to N$, then $d\phi$ induces an isomorphism
$N_{S/M} \cong S \times T_qN$, so $N_{S/M}$ is trivial.
\end{proposition}

\begin{proof}
The question is local, so let $E = U \times \mathbb{K}^r$, $F = U \times \mathbb{K}^s$ and $f$ a smooth family of matrices $A(p)$ of constant rank $k$. Near $p_0$ some $k \times k$ minor of $A$ is nonzero; after permuting bases,
$A = \begin{pmatrix}A_{11} & A_{12}\\ A_{21} & A_{22}\end{pmatrix}$ with $A_{11}$ invertible near $p_0$. The first $k$ columns of $A$ are linearly independent and, as the rank is $k$, span the image; so they form a frame of $\operatorname{im}f$,
which is therefore a subbundle. The vectors $\begin{pmatrix}-A_{11}^{-1}A_{12}w\\ w\end{pmatrix}$, $w \in \mathbb{K}^{r-k}$, lie in the kernel (the last $s - k$ rows are combinations of the first $k$), and they span it by dimension count; they
form a frame of $\ker f$ depending smoothly on $p$. Completing a frame of $\operatorname{im}f$ by constant vectors to a frame of $F$ gives a frame of $\operatorname{coker}f$. For the normal bundle apply this to the inclusion
$TS \to TM|_S$, of constant rank $\dim S$. If $S = \phi^{-1}(q)$ with $q$ regular, $d\phi \colon TM|_S \to S \times T_qN$ is surjective with kernel $TS$ (Theorem~\ref{smooth-manifolds:theorem-regular-value}), so it induces the stated isomorphism.
\end{proof}

\begin{counterexample}\label{vector-bundles:counterexample-rank-jump}
The rank must be constant. The morphism $f \colon \RR \times \RR \to \RR \times \RR$, $(x, t) \mapsto (x, xt)$, of trivial line bundles over $\RR$ has rank $1$ for $x \neq 0$ and $0$ at
$x = 0$. Its kernel is the zero vector over each $x \neq 0$ and the whole line over $x = 0$, and its image is the whole line over each $x \neq 0$ and the zero vector over $x = 0$.
In both cases the dimension of the fibres jumps, so neither is a subbundle.
\end{counterexample}

\begin{example}\label{vector-bundles:example-stably-trivial}
$TS^n$ is \emph{stably trivial}: $TS^n \oplus (S^n \times \RR) \cong S^n \times \RR^{n+1}$, because the normal bundle of $S^n = f^{-1}(1)$, $f(x) = |x|^2$, is trivial by Proposition~\ref{vector-bundles:proposition-constant-rank}, and
$TS^n \oplus N_{S^n/\RR^{n+1}} \cong T\RR^{n+1}|_{S^n}$ by Proposition~\ref{vector-bundles:proposition-metrics}(1). Concretely, $TS^n$ is the kernel of the constant rank morphism
$S^n \times \RR^{n+1} \to S^n \times \RR$, $(x, v) \mapsto (x, \langle x, v\rangle)$, and the normal line is spanned by the section $x \mapsto x$. Nevertheless $TS^2$ is not trivial,
since $S^2$ has no nowhere vanishing vector field (the hairy ball theorem, Corollary~\ref{singular-homology:corollary-sphere-applications}); so a bundle is not determined by its stable class.
\end{example}

\begin{example}\label{vector-bundles:example-euler-sequence}
On $\CC\PP^n$ there is an exact sequence of holomorphic vector bundles
$0 \to \mathcal{O} \to \mathcal{O}(1)^{\oplus(n+1)} \to T\CC\PP^n \to 0$ (the \emph{Euler sequence}), where $\mathcal{O}(1) = \mathcal{O}(-1)^\vee$ (Chapter~\ref{holomorphic-bundles}). Over the smooth category it
splits by Proposition~\ref{vector-bundles:proposition-metrics}, but it does not split holomorphically: holomorphic splittings are much rarer, which is why holomorphic bundles carry more information than
smooth ones.
\end{example}

\section{Line bundles and the first Chern class}\label{vector-bundles:section-line-bundles}

\begin{definition}\label{vector-bundles:definition-picard}
The \emph{Picard group} $\Pic^\infty(M)$ is the set of isomorphism classes of smooth complex line bundles with the tensor product; by Proposition~\ref{sheaves:proposition-modules} it is a group, with
inverse the dual.
\end{definition}

For line bundles the transition functions are nowhere vanishing functions $g_{ij} \colon U_i \cap U_j \to \CC^\times$, which commute, and the tensor product multiplies them. So
$\Pic^\infty(M)$ is the first \v{C}ech cohomology group of the sheaf of nowhere vanishing smooth functions, and it can be computed with the exponential function.

\begin{theorem}\label{vector-bundles:theorem-line-bundles}
Let $\underline{\CC}^\times_M$ be the sheaf of smooth nowhere vanishing complex functions. There are natural isomorphisms
\[
\Pic^\infty(M) \cong H^1(M, \underline{\CC}^\times_M) \xrightarrow{\ \sim\ } H^2(M; \ZZ),
\]
the second induced by the connecting homomorphism of the \emph{exponential sequence}
$0 \to \ZZ \to \underline{\CC}_M \xrightarrow{\exp(2\pi i\,\cdot)} \underline{\CC}^\times_M \to 0$. The resulting class $c_1(L) \in H^2(M; \ZZ)$ is the \emph{first Chern class}; it satisfies
$c_1(L \otimes L') = c_1(L) + c_1(L')$, $c_1(f^*L) = f^*c_1(L)$, and $L$ is trivial if and only if $c_1(L) = 0$.
\end{theorem}

\begin{proof}
Line bundles are classified by $\check{H}^1$ of the sheaf of smooth nonvanishing functions by Proposition~\ref{vector-bundles:proposition-cocycles}(3), and \v{C}ech and sheaf cohomology agree in degree one
(Remark~\ref{sheaf-cohomology:remark-cech-limit}); the group structures agree because the cocycle of a tensor product is the product of the cocycles. The exponential sequence is exact as a sequence of
sheaves: surjectivity holds because a nonvanishing function has a smooth logarithm on any set where it is defined and contractible, for instance a small ball, and the kernel of $\exp(2\pi i\,\cdot)$ is the
constant sheaf $\ZZ$ (Theorem~\ref{real-analysis:theorem-pi}). Its long exact sequence reads
\[
H^1(M, \underline{\CC}_M) \to H^1(M, \underline{\CC}^\times_M) \to H^2(M, \ZZ_M) \to H^2(M, \underline{\CC}_M),
\]
and the outer terms vanish because $\underline{\CC}_M$ is a sheaf of modules over $\mathcal{C}_M$, hence fine, hence acyclic (Corollary~\ref{sheaf-cohomology:corollary-fine-acyclic}); manifolds are
paracompact (Proposition~\ref{general-topology:proposition-exhaustion}). Finally $H^2(M, \ZZ_M) \cong H^2(M; \ZZ)$ by Theorem~\ref{sheaf-cohomology:theorem-singular-comparison}, since manifolds are locally
contractible. Additivity and naturality follow from the corresponding properties of the cocycles and of the connecting homomorphism, and triviality of $L$ means exactly that its class in
$H^1(M, \underline{\CC}^\times_M)$ vanishes.
\end{proof}

In particular every complex line bundle on a manifold with $H^2(M; \ZZ) = 0$ is trivial: on $\RR^n$, on a disc, on $S^1$, on $S^3$. On the sphere $S^2$ the theorem gives
$\Pic^\infty(S^2) \cong \ZZ$, and in this case the classification can be made completely explicit, by the winding number of a transition function.

For a continuous loop $c \colon [0, 1] \to \CC^\times$, $c(0) = c(1)$, choose a continuous logarithm $L$ along it, $e^{L(t)} = c(t)$; one exists because $\exp \colon \CC \to \CC^\times$ is a covering
map (local inverses are the branches of the logarithm on slit planes) and paths lift (Proposition~\ref{homotopy:proposition-homotopy-lifting}). The \emph{winding number}
$\mathrm{wind}(c) = (L(1) - L(0))/2\pi i$ is an integer, as $e^{L(1)} = e^{L(0)}$, and does not depend on $L$, since two lifts differ by a constant in $2\pi i\ZZ$. For a function
$g \colon \CC^\times \to \CC^\times$ we write $\mathrm{wind}(g)$ for the winding number of $t \mapsto g(e^{2\pi it})$ along the unit circle.

\begin{lemma}\label{vector-bundles:lemma-winding-number}
\begin{enumerate}
\item $\mathrm{wind}(c_1c_2) = \mathrm{wind}(c_1) + \mathrm{wind}(c_2)$, and $\mathrm{wind}(z^k) = k$.
\item Loops that are homotopic through loops in $\CC^\times$ have the same winding number. In particular $\mathrm{wind}(g) = 0$ if $g$ extends to a continuous nowhere vanishing
function on the closed unit disc, or on $\{|z| \geq 1\} \cup \{\infty\} \subseteq \CC\PP^1$.
\item A continuous $g \colon \CC^\times \to \CC^\times$ with $\mathrm{wind}(g) = 0$ has a continuous logarithm on $\CC^\times$, smooth if $g$ is smooth.
\end{enumerate}
\end{lemma}

\begin{proof}
(1) If $L_1$, $L_2$ are logarithms along $c_1$, $c_2$, then $L_1 + L_2$ is one along $c_1c_2$. For $z^k$ along $e^{2\pi it}$ take $L(t) = 2\pi ikt$. (2) Let $H(s, t)$ be a homotopy of loops.
Lifting $H$ through $\exp$ (Proposition~\ref{homotopy:proposition-homotopy-lifting}) gives $\tilde H$ with $e^{\tilde H} = H$, and $s \mapsto (\tilde H(s, 1) - \tilde H(s, 0))/2\pi i$ is a
continuous integer-valued function, hence constant. If $g$ extends to the closed unit disc, the loops $t \mapsto g(\rho e^{2\pi it})$, $0 \leq \rho \leq 1$, form a homotopy through loops in
$\CC^\times$ from the constant loop $g(0)$ to the given one; if it extends to $\{|z| \geq 1\} \cup \{\infty\}$, use the loops $t \mapsto g(\sigma^{-1}e^{2\pi it})$, $0 < \sigma \leq 1$, with the
constant loop $g(\infty)$ at $\sigma = 0$.
(3) The function $G(r, t) = g(re^{2\pi it})$ on $(0, \infty) \times \RR$, which is simply connected, lifts through $\exp$ to a continuous $\Lambda$ with $e^\Lambda = G$
(Proposition~\ref{homotopy:proposition-lifting-criterion}). The difference $\Lambda(r, t + 1) - \Lambda(r, t)$ lies in $2\pi i\ZZ$ and depends continuously on $(r, t)$, so it is a constant, equal to
$2\pi i$ times the winding number of the circle of radius $1$, which is $0$. So $\Lambda(r, t)$ depends only on $re^{2\pi it}$, and defines a continuous logarithm of $g$ on $\CC^\times$; it is
smooth when $g$ is, because $\exp$ is a local diffeomorphism.
\end{proof}

\begin{proposition}\label{vector-bundles:proposition-line-bundles-sphere}
Let $L$ be a complex line bundle on $\CC\PP^1 = U_0 \cup U_1$, with $U_0 = \{[1 : z]\} \cong \CC$ and $U_1 = \{[w : 1]\}$, $w = 1/z$ on the overlap. Then $L$ has frames $s_0$ over
$U_0$ and $s_1$ over $U_1$, and on $U_0 \cap U_1 \cong \CC^\times$ (coordinate $z$) we can write $s_1 = g\,s_0$ with $g \colon \CC^\times \to \CC^\times$. The integer
\[
\deg L = \mathrm{wind}(g)
\]
does not depend on the frames, $\deg(L \otimes L^{\prime}) = \deg L + \deg L^{\prime}$, every integer is the degree of some line bundle, and $L \cong L^{\prime}$ if and only if
$\deg L = \deg L^{\prime}$. So $\deg \colon \Pic^\infty(\CC\PP^1) \to \ZZ$ is an isomorphism.
\end{proposition}

\begin{proof}
The frames exist because $U_0$ and $U_1$ are diffeomorphic to $\RR^2$, over which every complex line bundle is trivial, as $H^2(\RR^2; \ZZ) = 0$
(Theorem~\ref{vector-bundles:theorem-line-bundles}). Other frames are $s^{\prime}_0 = h_0s_0$ and $s^{\prime}_1 = h_1s_1$ with $h_i \colon U_i \to \CC^\times$, and then
$g^{\prime} = h_1gh_0^{-1}$; by Lemma~\ref{vector-bundles:lemma-winding-number}(1),(2), $\mathrm{wind}(g^{\prime}) = \mathrm{wind}(g) + \mathrm{wind}(h_1) - \mathrm{wind}(h_0) = \mathrm{wind}(g)$, as $h_0$ is
defined on all of $U_0 = \CC$ and $h_1$ on $U_1 = \CC\PP^1 \setminus \{0\}$, which contains $\{|z| \geq 1\} \cup \{\infty\}$. For the tensor product, $s_1 \otimes s^{\prime}_1 = gg^{\prime}\,s_0 \otimes s^{\prime}_0$, so
degrees add. The bundle with transition function $g = z^k$, obtained by Proposition~\ref{vector-bundles:proposition-cocycles}(2), has degree $k$. If $\deg L = \deg L^{\prime}$, with transition
functions $g$ and $g^{\prime}$, then $g^{\prime}/g$ has winding number $0$, hence a smooth logarithm $f$ on $\CC^\times$ (Lemma~\ref{vector-bundles:lemma-winding-number}(3)). Choose a partition of unity
$\psi_0 + \psi_1 = 1$ subordinate to $\{U_0, U_1\}$ and put $f_0 = -\psi_1f$ on $U_0$ and $f_1 = \psi_0f$ on $U_1$, each extended by $0$ where the cut-off vanishes; then $f_1 - f_0 = f$ on the
overlap, and $h_i = e^{f_i}$ satisfy $h_1gh_0^{-1} = ge^{f} = g^{\prime}$. So after changing the frames of $L$ by $h_0$ and $h_1$, the two bundles have the same transition function, and they are
isomorphic by Proposition~\ref{vector-bundles:proposition-cocycles}(3).
\end{proof}

\begin{example}\label{vector-bundles:example-degrees}
\begin{enumerate}
\item The tautological bundle has $s_1 = z^{-1}s_0$ (Example~\ref{vector-bundles:example-cocycles}(4)), so $\deg\mathcal{O}(-1) = -1$, and $\deg\mathcal{O}(k) = k$.
\item The tangent bundle $T\CC\PP^1$, a complex line bundle with the frames $s_0 = \partial/\partial z$ and $s_1 = \partial/\partial w$: from $z = 1/w$,
$\partial/\partial w = (dz/dw)\,\partial/\partial z = -w^{-2}\,\partial/\partial z = -z^2\,\partial/\partial z$, so $g = -z^2$ and $\deg T\CC\PP^1 = 2$. Hence
$T\CC\PP^1 \cong \mathcal{O}(2)$, and $TS^2$ is not trivial, which gives another proof of the hairy ball theorem for $S^2$. The cotangent bundle has degree $-2$.
\item The degree is the first Chern class: under $H^2(\CC\PP^1; \ZZ) \cong \ZZ$ given by the complex orientation, $c_1(L) = \deg L$. We do not verify this here by computing the
connecting homomorphism; it follows from Chern--Weil theory (Chapter~\ref{characteristic-classes}), and in Chapter~\ref{connections} the curvature of a connection on
$\mathcal{O}(-1)$ is computed to integrate to $-1$. The degree $2$ of $T\CC\PP^1$ is the Euler characteristic of $S^2$, an instance of the Gauss--Bonnet theorem.
\end{enumerate}
\end{example}

\begin{example}\label{vector-bundles:example-line-bundles}
$H^2(S^2; \ZZ) \cong \ZZ$, so the complex line bundles on $S^2 = \CC\PP^1$ are classified by an integer, the degree; the tautological bundle $\mathcal{O}(-1)$ has $c_1 = -1$ with the standard
orientation. On $\CC\PP^n$, $H^2 \cong \ZZ$ is generated by $c_1(\mathcal{O}(1))$. Every complex line bundle on a surface of genus $g$ is classified by its degree in $H^2 \cong \ZZ$, while on a simply
connected manifold with $H^2 = 0$, such as $S^3$ or $S^4$, every complex line bundle is trivial. Real line bundles are classified in the same way by
$H^1(M; \ZZ/2\ZZ)$, using the sheaf of nonvanishing real functions and the sequence $0 \to \ZZ/2 \to \underline{\RR}^\times \to \underline{\RR}^\times_{>0} \to 0$; the resulting class is the first
Stiefel--Whitney class $w_1(L)$, which vanishes if and only if $L$ is orientable.
\end{example}

\begin{remark}\label{vector-bundles:remark-higher-chern}
For bundles of higher rank the classification is not by a single class: isomorphism classes of rank $r$ complex bundles on $M$ correspond to homotopy classes of maps $M \to \mathrm{BU}(r)$, and the Chern
classes $c_k(E) \in H^{2k}(M; \ZZ)$ are pulled back from universal classes on the classifying space. They are constructed and computed in Chapter~\ref{characteristic-classes}; for line bundles $c_1$ agrees
with the class above.
\end{remark}

\section{Exercises}\label{vector-bundles:section-exercises}

\begin{exercise}\label{vector-bundles:exercise-mobius}
Show that the M\"obius bundle $L \to S^1$ is nontrivial, that $L \otimes L$ is trivial, and that $L$ has no nowhere vanishing section.
\end{exercise}

\begin{solution}
A nowhere vanishing section would give a trivialisation; on the other hand a section is a continuous function $f \colon [0, 1] \to \RR$ with $f(1) = -f(0)$, which must vanish somewhere by the intermediate
value theorem. The transition function of $L \otimes L$ is the square of that of $L$, namely $1$.
\end{solution}

\begin{exercise}\label{vector-bundles:exercise-tangent-sphere-nontrivial}
Show that $TS^2$ is not trivial, using that a trivialisation would give a nowhere vanishing vector field (Corollary~\ref{singular-homology:corollary-sphere-applications}).
\end{exercise}

\begin{solution}
A trivialisation $TS^2 \cong S^2 \times \RR^2$ gives the global frame $\Phi^{-1}(p, e_1)$, $\Phi^{-1}(p, e_2)$, and in particular the nowhere vanishing vector field $\Phi^{-1}(p, e_1)$,
which does not exist. For another proof, see Example~\ref{vector-bundles:example-degrees}(2).
\end{solution}

\begin{exercise}\label{vector-bundles:exercise-tautological}
Show that the tautological line bundle on $\CC\PP^n$ is $\mathcal{O}(-1) = \{(\ell, v) : v \in \ell\} \subseteq \CC\PP^n \times \CC^{n+1}$, that it is a submanifold and a line bundle, and compute its transition
functions for the standard cover.
\end{exercise}

\begin{solution}
On $U_i = \{x_i \neq 0\}$ the map $\ell \mapsto (\ell, (x_0/x_i, \ldots, 1, \ldots, x_n/x_i))$ is a nowhere vanishing section $s_i$, giving a trivialisation; then $s_j = (x_i/x_j)s_i$, so
$g_{ij} = x_i/x_j$ as functions on $U_i \cap U_j$.
\end{solution}

\begin{exercise}\label{vector-bundles:exercise-normal-bundle}
Let $S \subseteq M$ be a submanifold with normal bundle $N = (TM|_S)/TS$. Show that if $S = f^{-1}(q)$ for a submersion $f$, then $N \cong S \times \RR^m$ is trivial, and that the normal bundle of
$\RR\PP^{n-1} \subseteq \RR\PP^n$ is the tautological bundle.
\end{exercise}

\begin{solution}
The first statement is Proposition~\ref{vector-bundles:proposition-constant-rank}. For the second, $\RR\PP^{n-1} = \{x_n = 0\}$. A tangent vector to $\RR\PP^n$ at a line $\ell$ is the same as a
linear map $\ell \to \ell^\perp$ (the derivative of a curve of lines $\ell_t$, measured by how a vector of $\ell$ moves orthogonally), so $T_\ell\RR\PP^n = \Hom(\ell, \ell^\perp)$. For
$\ell \subseteq \RR^n = \{x_n = 0\}$, the tangent space of $\RR\PP^{n-1}$ is $\Hom(\ell, \ell^\perp \cap \RR^n)$, and the quotient is $\Hom(\ell, \RR e_n) \cong \ell^\vee$. So the normal bundle is the
dual of the tautological bundle, which is isomorphic to it by Proposition~\ref{vector-bundles:proposition-metrics}(2).
\end{solution}

\begin{exercise}\label{vector-bundles:exercise-inversion-jacobian}
Compute the Jacobian matrix of the inversion $\tau(y) = y/|y|^2$ of $\RR^n \setminus \{0\}$ and verify the formula of Example~\ref{vector-bundles:example-cocycles}(2), including
$\det D\tau(y) = -|y|^{-2n}$. For $n = 2$, write $y = x_1 + ix_2$ and check that $\tau(y) = 1/\bar y$.
\end{exercise}

\begin{solution}
$\partial_j(y_i/|y|^2) = \delta_{ij}/|y|^2 - 2y_iy_j/|y|^4$, which is the formula. The matrix $I - 2uu^T$ with $u = y/|y|$ is the reflection in $u^\perp$, of determinant $-1$, so
$\det D\tau = (|y|^{-2})^n \cdot (-1)$. For $n = 2$, $1/\bar y = y/(y\bar y) = y/|y|^2$.
\end{solution}

\begin{exercise}\label{vector-bundles:exercise-pullback-degree}
Let $f \colon \CC\PP^1 \to \CC\PP^1$ be $[x_0 : x_1] \mapsto [x_0^k : x_1^k]$ for $k \geq 1$. Show that $f^*\mathcal{O}(-1) \cong \mathcal{O}(-k)$, by computing the transition function of the
pulled back bundle.
\end{exercise}

\begin{solution}
In the coordinate $z$ of $U_0$, $f(z) = z^k$, and $f$ maps $U_0$ to $U_0$ and $U_1$ to $U_1$. The pulled back frames $s_0 \circ f$ and $s_1 \circ f$ satisfy
$s_1 \circ f = (z^k)^{-1}\,s_0 \circ f$, by Example~\ref{vector-bundles:example-cocycles}(4) evaluated at $f(z)$. So the transition function is $z^{-k}$, of winding number $-k$, and
$f^*\mathcal{O}(-1) \cong \mathcal{O}(-k)$ by Proposition~\ref{vector-bundles:proposition-line-bundles-sphere}. The degree is multiplied by the degree $k$ of the map.
\end{solution}

\begin{exercise}\label{vector-bundles:exercise-real-line-bundles}
Show that for a real line bundle $L$ the following are equivalent: $L$ is trivial; $L$ has a nowhere vanishing section; the transition functions of some trivialising cover can be
chosen positive. Deduce that the tautological bundle of $\RR\PP^n$ is not trivial for $n \geq 1$, and that every real line bundle over a simply connected manifold is trivial (using
the classification by $H^1(M; \ZZ/2)$ of Example~\ref{vector-bundles:example-line-bundles}).
\end{exercise}

\begin{solution}
A nowhere vanishing section $s$ gives the trivialisation $(p, t) \mapsto ts(p)$; conversely $p \mapsto \Phi^{-1}(p, 1)$. With positive transition functions, choose a partition of unity
and glue the local sections $\Phi_i^{-1}(p, 1)$: $s = \sum_i\psi_i\Phi_i^{-1}(\cdot, 1)$ is nowhere zero, because in any trivialisation the summands are positive multiples of each other.
Conversely, with a nowhere vanishing section, rescale the trivialisations so that it becomes the constant $1$; then all transition functions are $1$. For $\RR\PP^n$ the tautological bundle
restricts on $\RR\PP^1$ to the M\"obius bundle (Counterexample~\ref{vector-bundles:counterexample-mobius}). On a simply connected $M$, $H^1(M; \ZZ/2) = \Hom(\pi_1(M), \ZZ/2) = 0$.
\end{solution}

\begin{exercise}\label{vector-bundles:exercise-hom-cocycle}
Show directly that $\Hom(E, F)$, defined as the bundle with fibres $\Hom(E_p, F_p)$ and transition functions $A \mapsto h_{ij}Ag_{ij}^{-1}$, is isomorphic to $E^\vee \otimes F$, and that its sections
are exactly the morphisms of vector bundles $E \to F$.
\end{exercise}

\begin{solution}
Fibrewise $E_p^\vee \otimes F_p \to \Hom(E_p, F_p)$, $\phi \otimes w \mapsto (v \mapsto \phi(v)w)$, is an isomorphism (Proposition~\ref{multilinear-algebra:proposition-hom-tensor}); in trivialisations
it is a fixed linear isomorphism $\mathbb{K}^{r\vee} \otimes \mathbb{K}^s \to M_{s \times r}(\mathbb{K})$ intertwining the transition functions, since $(g^T)^{-1}\xi \otimes hw \mapsto hw\,\xi^Tg^{-1}$. So
it is an isomorphism of bundles (Lemma~\ref{vector-bundles:lemma-bijective-morphism}). A section is a smoothly varying family of linear maps $E_p \to F_p$, which is a morphism of bundles:
in trivialisations it is a smooth matrix-valued function.
\end{solution}

\begin{exercise}\label{vector-bundles:exercise-sum-mobius-explicit}
Let $E$ be the M\"obius bundle. Show that $E \oplus (S^1 \times \RR)$ is not trivial but that $E \oplus E$ is, by describing sections as functions on $\RR$ with the appropriate periodicity.
\end{exercise}

\begin{solution}
Sections of $E \oplus (S^1 \times \RR)$ are pairs $(f, g)$ with $f(x + 1) = -f(x)$ and $g(x + 1) = g(x)$. If it were trivial, its determinant line bundle $\det(E \oplus \RR) \cong E$
(Proposition~\ref{vector-bundles:proposition-operations}(3)) would be trivial. For $E \oplus E$, the sections $(\cos\pi x, \sin\pi x)$ and $(-\sin\pi x, \cos\pi x)$ are anti-periodic, hence
sections, and linearly independent everywhere (Example~\ref{vector-bundles:example-mobius-square}).
\end{solution}
